\section{Future Work}
\label{sec:future}

Our study establishes that a typed refutation store removes verification cost
and that its prompt-side half does not work at $7$B scale. Two programmes
follow: a set of baseline arms that would place \CEGMem\ against the mechanisms
other lines of work actually use, and a set of extensions the negative result
points at. We describe the first in enough detail to be pre-specified, because
we think the design rules matter more than the outcomes we would report.

\subsection{Baseline arms, and the rules for running them}
\label{sec:future-arms}

The comparison we want is not against published \emph{numbers}---other systems
run other models on other benchmarks, and a reviewer is right to reject that as
unfair. It is against published \emph{mechanisms}, reimplemented inside this
harness so that model, corpus, response cache and common-random-number pairing
are held fixed and only the mechanism varies. \Cref{tab:futurearms} lists the
five arms; three are implemented and pass a round-1 pairing test, and the
budget for them is the reason they are future rather than present work.

\begin{table*}[t]
\centering
\small
\caption{Planned baseline arms (\cref{sec:future-arms}). Each reimplements a
published \emph{mechanism} inside this harness---same model, corpus, response
cache and common-random-number pairing---so only the mechanism varies. Three
are implemented and pass the round-1 pairing test; none has been run at scale,
which is why they appear here rather than in \cref{sec:results}. The last
column states the prediction we would be pre-specifying, including where we
expect to be wrong.}
\label{tab:futurearms}
\setlength{\tabcolsep}{4.5pt}
\renewcommand{\arraystretch}{1.12}
\begin{tabular}{@{}L{1.6cm}L{2.5cm}L{5.0cm}L{1.5cm}L{4.4cm}@{}}
\toprule
\textbf{Arm} & \textbf{Mechanism (source)} & \textbf{Configuration in this
harness} & \textbf{Cells} & \textbf{Pre-registered prediction} \\
\midrule

conversation &
ChatRepair's growing dialogue~\cite{xia2024chatrepair} &
guard off, steer off; the prompt carries every failed patch \emph{in full} plus
the counterexample that killed it; conversation restarts after $5$ consecutive
failures or ${\approx}6$k transcript tokens; restarts logged &
$297$ \sig{(corpus, seeds 1--3)} &
Outcome close to no memory, tokens far above it. Unlike our \armsteer, it shows
the \emph{patches}, so it may reduce the $21\%$ verbatim-duplicate rate---the
one place we expect a real effect and would be reporting a surprise. \\

self-test &
CodeT's generated-test filter~\cite{chen2023codet} &
guard off, steer off; one cached model call generates test cases per fault;
every proposal must pass them \emph{before} the oracle is paid &
$90$ \sig{(sweep)} &
The only baseline that attacks oracle cost the way the guard does, but with
a-priori model knowledge instead of accumulated refutations. We would report
filter precision against the oracle's verdict, since a wrong generated test
blocks a \emph{correct} patch. \\

self-test $+$ guard &
composition &
as above with the typed guard on &
$90$ &
Self-tests catch a-priori failures, the guard catches repeats; if the two are
complementary this is the lowest oracle cost in the study. \\

random-skip &
control for ``execution applied
indiscriminately''~\cite{lin2026run} &
guard off; before each oracle call, skip with $p = 0.37$---the typed guard's
measured block share---seeded per cell; a skipped round costs budget exactly as
a blocked one does &
$90$ &
Repair should \emph{fall}, because skipping at random discards correct patches
in proportion. This is the \cref{sec:results-typing} falsifier one level up: it
isolates ``knowing \emph{which} call to skip'' as the mechanism. \\

reflection &
Reflexion / Self-Debug verbal
memory~\cite{shinn2023reflexion,chen2024selfdebug} &
guard off, steer off; after each refutation one extra call writes a short
reflection (${\le}120$ tokens, last $3$ retained) inserted into the next
prompt; separate nonce so the proposal sequence is unperturbed &
$90$ \sig{(appendix)} &
Null expected at $7$B, on the same reasoning as \cref{sec:discussion}. Its
value is defensive: it answers ``raw evidence failed, but would distilled
reflection have worked?'' It is also the most expensive arm per cell, since a
reflection never hits the cache. \\

\bottomrule
\end{tabular}
\end{table*}


Four rules govern the family, and we state them before running it rather than
after.

\emph{Protocol invariance.} No arm may change the model, temperature, output
budget, test tree or sandbox timeout, and the proposal nonce keeps its
semantics, so round~1 remains the identical draw in every arm---old and new.
Each arm ships a test asserting that identity, because an arm that breaks the
pairing silently converts every comparison in this paper into an unpaired one.

\emph{The primary metric stays the budget curve.} Every arm here lacks a guard
and therefore loses a total-oracle-call comparison \emph{by construction}. We
pre-specify repair rate at a matched oracle budget as primary and report total
calls as descriptive, for the same reason we did in \cref{sec:design}: a
tautology is not a result, and a reviewer should not have to point that out.

\emph{Multiple comparisons declared in advance.} The family is
$\{$five arms$\} \times \{$repair at $b{=}8$, oracle calls$\}$ and its
$p$-values will be Benjamini--Hochberg adjusted~\cite{benjamini1995fdr}
over the whole family, not per arm.

\emph{A fidelity table per arm.} Each reimplementation states which features of
the original it ports, which it omits, and why---for instance, our
conversation arm ports history plus the restart policy but not the request for
variations around a plausible patch. A mechanism port with an undocumented
omission is a straw man, and the omission is where a reviewer will look first.

\subsection{Two comparisons that need no new runs}
Two questions in the literature can be answered from the frozen log.
\emph{One-shot versus iterative repair}: cutting the steer-only and no-memory
logs at round~2 recovers exactly the question
Olausson et al.~\cite{olausson2024selfrepair} ask---is the second attempt, given feedback,
better than a fresh draw?---paired per cell, at zero cost.
\emph{Sample-and-verify}: because the no-memory arm's draws are i.i.d.\ samples
under a fixed nonce, its budget curve \emph{is} the pass@$k$-with-verification
baseline, and \cref{fig:budget} already plots it. We note both here rather than
in \cref{sec:results} because neither is a new measurement; they are
re-readings of the same log against two other papers' framings.

\subsection{Directions the negative result points at}
\emph{Adaptive steering.} The prompt blocks are worth their $+128\%$ tokens
only where the baseline is near zero (\cref{sec:discussion}). An arm that
enables them after $m$ consecutive failures, with $m$ set from the per-band
anchoring and success measurements, is the cheapest test of whether the
dead-band gain survives being targeted.

\emph{A guard for partial oracles.} Our guard re-executes a stored
counterexample, which presumes a counterexample can be stored and replayed
cheaply. At repository scale, the analogue is a stored \emph{failing test plus
its environment}, and the interesting question is whether a type map over stack
traces and changed files concentrates refutations better than $\loc$ does. This is where the
mechanism would meet SWE-bench-scale work~\cite{jimenez2024swebench,%
yang2024sweagent,xia2025agentless}, and where the guard's arithmetic is
most favourable, because each avoided validation is worth minutes rather than
seconds.

\emph{Composition with faster validation.} A guard removes calls and
UniAPR-style validation makes the survivors cheap~\cite{chen2021uniapr,%
xiao2024expressapr}. The two have never been measured together, and the
product of \fold{5} and \fold{100} is the interesting number.

\subsection{Directions we considered and rejected}
We record three, with reasons, because each looks attractive and each would
break something this study depends on.
\emph{Cross-task retrieval}---sharing memory across faults---breaks the
independence our corpus is built on and destroys the per-task pairing that
every $p$-value here rests on; with $99$ unrelated single-file faults the
expected transfer is also near zero.
\emph{Greedy decoding} ($T = 0$) would invalidate the response cache and, more
fundamentally, make redraw-after-block vacuous (\cref{sec:design}); the
question it is meant to answer---how much of the loop is repeated
sampling---is already answered by the measured $\pihat = 0.207$ and the
pass@$1$ of $0.214$ against pass@$20$ of $0.703$.
\emph{Seed-pooled memory}---accumulating memory across seeds of the same
task---would raise the block rate but collapse the replication axis the
statistics are computed over, trading the study's inferential structure for a
larger single number.
